% ============================================================================
\chapter{Ступень VII. Кремниевый чип и инженерные установки своими руками}
\label{ch:course-stage7-engineering}
% ============================================================================

\section{Урок 22. Как устроен наш 9-битный бесфлаговый чип (\texttt{cuboctahedron\_carrier\_asic.v}) и аналоговая 3D LC-решётка (\texttt{cuboctahedron\_12beam\_spice.cir})}

Теория, которую нельзя спаять и проверить осциллографом, --- это только половина квеста. Поэтому вся наша математика из Ступеней~I--VI переведена в два открытых аппаратных ядра:

\begin{quote}
\small\textbf{Честная инженерная оговорка и приглашение к стендовой проверке.}
Все представленные в этой главе цифровые схемы (\texttt{Verilog}), аналоговые контуры (\texttt{SPICE}), чертежи плат (\texttt{Gerber}) и спецификации деталей прошли логическую и схемотехническую симуляцию (\texttt{Icarus Verilog}, \texttt{LTspice}). Но симуляция --- это ещё не физический кремний и не паяльник на столе: \textbf{в реальных схемах, разводке и протоколах могут быть ошибки, неучтённые паразитные ёмкости или допуски; что-то может не заработать с первой сборки и потребует доработки, дополнения и независимой перепроверки инженерами и экспериментаторами}. Все чертежи и коды полностью открыты как Википедия --- мы приглашаем инженеров со всего мира проверять, исправлять и доводить их на реальных стендах вместе.
\end{quote}

\begin{enumerate}[label=\textbf{\arabic*.},leftmargin=*]
\item \textbf{Цифровой RTL-чип на ПЛИС (\texttt{hardware/cuboctahedron\_carrier\_asic.v}, 16 модулей, 0 ошибок в \texttt{iverilog}):}
  \begin{itemize}[noitemsep]
  \item Единственный параметр всего чипа --- размерность пространства \texttt{parameter integer D = 3}; даже единица вычисляется как \texttt{localparam integer MU\_VAC = (D * D) - (1 << D);}.
  \item В вычислительном ядре \textbf{0 операторов \texttt{if}, 0 тернарных мультиплексоров \texttt{?:}, 0 таблиц памяти синусов и 0 битовых сдвигов на 9-битных кольцевых шинах}: ротатор (\texttt{Module~1}) работает на бабочке Кода Грея $(128 - m_7, m_7)$, безмножительное АЛУ (\texttt{Module~1B}) перемножает знаковые числа через первообразную Верле $\mathcal{F}_S$, расщепитель (\texttt{Module~2}) берёт волну с 9 проводов суммы и корпускулу с проводов переноса, двухтактовая клетка Верле (\texttt{Module~3}) обращает время без стека истории, а 12-соседнее ядро КА (\texttt{Modules~3B, 3C}) параллельно опрашивает 12 аппаратных ротаторов по векторам кубооктаэдра~$O_2$.
  \item Топ-модуль SoC (\texttt{Module~12}) и физический UART-десериализатор платы \textbf{Sipeed Tang Nano 9K} (\texttt{Module~16}, \texttt{hardware/tang\_nano\_9k.cst}, 16-байтный бинарный кадр \texttt{[0xA5, op, arg0..arg2, xor\_crc, 0x5A]} на $115\,200$ бод) принимают команды с ноутбука/ИИ по USB-UART или Wi-Fi и поддерживают \textbf{горячий патч на лету} (опкоды \texttt{4'h6}, \texttt{4'h7} и \texttt{4'h8}), позволяющий без перепрошивки чипа задать любую индивидуальную фазу $0\dots 511$ на каждой из 12 вершин кубооктаэдра, любую маску активных лучей \texttt{beam\_oe[11:0]} и время выдержки каждой стадии синтезатора \texttt{stage\_hold\_ticks}.
  \item \textbf{Сам атом внутри камеры как ведущий сервер точного времени и полное устранение проблемы неизвестных координат (опкод \texttt{4'h8}):}
    как попасть 12-лучевым фазовым полем точно в атом внутри рабочей камеры, если из-за барьера Демона Лапласа ($64\times 81 = 72^2 = 5\,184$) мы не можем заранее измерить его микрокоординаты $(x,y,z)$ и начальную микрофазу $\varphi_{\mathrm{atom}}\in\mathbb{Z}_{512}$? Секрет в том, что машине \textbf{не нужен внешний независимый эталон времени}: сам атом в камере служит ей ведущим сервером времени (Atomic Self-Clocked PLL)! Нечётная проекция ротатора $e_\varphi = \mathcal{S}_3(\varphi_{\mathrm{atom}} - \varphi_{\mathrm{machine}})$ в мёртвой зоне Гейзенберга ($\hat{n}_E = 0$) работает как безмножительный фазовый детектор ($\mathcal{S}_3(0) = 0,\ \mathcal{C}_3(0) = +128$), захватывая фазу атома, после чего драйвер излучает обратно сопряжённую волну $\varphi_{\mathrm{drive}} = (\varphi_{\mathrm{atom}} + \Delta\varphi_{\mathrm{mode}})\bmod 512$. На самом атоме неизвестная фаза и путевая координата $\varphi_{\mathrm{atom}}$ \textbf{телескопически сокращаются в точный ноль на всех 512 состояниях кольца}:
    \[
      \bigl(\varphi_{\mathrm{drive}} - \varphi_{\mathrm{atom}}\bigr)\bmod 512
      \;=\;
      \bigl(\varphi_{\mathrm{atom}} + \Delta\varphi_{\mathrm{mode}} - \varphi_{\mathrm{atom}}\bigr)\bmod 512
      \;\equiv\; \mathbf{\Delta\varphi_{\mathrm{mode}}}
      \qquad\Longrightarrow\qquad
      D\bigl(\Delta\varphi_{\mathrm{mode}}=139\bigr) \;\equiv\; \mathbf{8\,192}
      \;\;\forall\,\varphi_{\mathrm{atom}}\in\mathbb{Z}_{512}!
    \]
  \item \textbf{4D-навигация по мировому конусу и обращение времени целого макрообъекта без вакуума:}
    более того, поскольку наша Вселенная на 100\% детерминирована от сида $\mathrm{IV}(3) = 263 = 407_8$, любой попавший в камеру физический атом уже несёт в себе фоновый градиент $\varphi_{\mathrm{bg}}(x) = (x\cdot 41)\bmod 512$ ($x\in\{0,\dots,12\}$) и скорость параболы Верле $\Delta_t\varphi = (t\cdot 41)\bmod 512$. Так как $12\times 41 = 492 < 512$ не перехлёстывает через кольцо, а обратным к $41$ по модулю $512$ является ключ Безу $\mathbf{25 = d_{\mathrm{out}}(E)^2}$ ($25\times 41 = 1025 = 2\times 512 + 1$), одно 9-битное умножение на $25\bmod 512$ мгновенно расшифровывает \textbf{все $13/13$ пространственных координат $x\in[0..12]$ и все $512/512$ тактов времени $t\in[0..511]$}:
    \[
      \bigl(25\cdot \varphi_{\mathrm{bg}}(x)\bigr)\bmod 512 \;\equiv\; \mathbf{x},
      \qquad
      \bigl(25\cdot \Delta_t\varphi\bigr)\bmod 512 \;\equiv\; \mathbf{t}!
    \]
    При этом не нужны ни глубокий вакуум, ни одноатомная ловушка: включение синфазного стазиса \texttt{Mode~0} ($D_{\mathrm{stasis}} = 6\,823\,936$) гасит тепловой разброс между $10^{22}$ атомами камеры в ноль мёртвой зоной Гейзенберга, сцепляя их по теореме о гомоморфизме среды в единый когерентный макро-ансамбль. А чтобы сдвинуть назад или вперёд по мировой линии Верле весь макрообъект целиком без разрыва связей на границе с комнатой, вокруг всей капсулы замыкается 12-лучевой марковский купол стазиса ($D_{\mathrm{dome}} = 833\,\varepsilon = 6\,823\,936 \implies \Phi_{\mathrm{ext\to int}} \equiv 0$), который полностью отсекает внешнюю стрелу времени комнаты и позволяет всему внутреннему объёму синхронно шагать по обратимому закону Верле!
  \item \textbf{Разрешение парадокса Хокинга: почему на вечеринку 2009~года не пришли «туристы из будущего»:}
    три строгих закона носителя объясняют это без единого противоречия:
    (а)~\emph{Закон сохранения материи A3 ($\sum \mu_x^2 = \mathrm{const}$):} нельзя «размножить» атомы ($1\to 2$) и материализовать в 2009~году вторую копию тела ($\mu=64$) рядом с уже существовавшими в 2009~году атомами ($64^2+64^2 = 8\,192 \ne 4\,097 = 64^2+1^2$);
    (б)~\emph{Внутреннее обращение времени кокона против глобальной ленты:} обратный шаг Верле $g_{\mathrm{int}}^{-N}$ внутри отрезанного стазис-купола ($D_{\mathrm{dome}} = 6\,823\,936$) отматывает в прошлое \textbf{собственное состояние самой капсулы} $\Psi_{\mathrm{cocoon}}(\tau)\mapsto\Psi_{\mathrm{cocoon}}(\tau-N)$ (позволяя откатить назад травму, старение или реакцию внутри купола, либо через стазис $g^0$ перешагнуть в будущее с нулевым собственным временем $\Delta\tau_{\mathrm{proper}}=0$), что исключает парадокс дедушки;
    (в)~\emph{Горизонт включения приёмника ($\tau \ge \tau_{\mathrm{boot}}$) и бесфотонная невидимость Канала~1 ($\hat{n}_E = 0$):} приём фазового пакета $(b + \sigma_A)\bmod 512 = \psi_{\mathrm{src}}$ возможен только там и тогда, где уже собрана и включена приёмная 12-лучевая станция с парным синглетом Паули $b = a + 256$ (которой в 2009~году у Хокинга в комнате не было!), а любое невозмущающее наблюдение по Каналу~1 в мёртвой зоне Гейзенберга $|\Delta\varphi|\le 40 < 41$ излучает $\hat{n}_E \equiv \mathbf{0\text{ фотонов}}$ и невидимо для обычных камер!
  \item \textbf{Универсальный Хроновизор и замкнутая парабола Верле за $O(1)$ шагов:}
    раз материальный прыжок в прошлое запрещён законом сохранения массы~A3, как \emph{увидеть} прошлое на экране или в голографическом куполе? Поскольку вся Вселенная детерминированно выросла из известного 9-битного слова $\mathrm{IV}(3) = 263 = 407_8$, нам не нужно измерять всю Вселенную сегодня (обход Демона Лапласа!) --- достаточно одного физического атома в комнате, который через ключ Безу $25\times 41 \equiv 1\pmod{512}$ даёт точный якорь наших текущих координат $(x_{\mathrm{now}}, \tau_{\mathrm{now}})$ на глобальной 4D-карте от $\mathrm{IV}=263$:
    \begin{itemize}[noitemsep]
    \item \emph{Режим~А (Генеалогический луч):} сам ребёнок в момент зачатия является физическим слиянием двух родительских мировых линий по теореме о гомоморфизме подсистем $\mathcal{E}_{\mathrm{child}}(\tau_{\mathrm{meet}}) = \mathcal{E}_{\mathrm{mom}}(\tau_{\mathrm{meet}}) \++ \mathcal{E}_{\mathrm{dad}}(\tau_{\mathrm{meet}})$. Поэтому обратный шаг Верле от атомов/ДНК самого человека ведёт траекторию назад по его мировой линии как по рельсам прямо в точку встречи родителей $(\mathbf{r}_{\mathrm{meet}}, \tau_{\mathrm{meet}})$!
    \item \emph{Режим~Б (Прыжок за $O(1)$ в любую произвольную точку пространства-времени $(\mathbf{r}_{\mathrm{target}}, \tau_{\mathrm{target}})$):} по теореме о замкнутой параболе Верле (так как вторая разность $\Delta^2\varphi(t) \equiv \Phi$ постоянна на стационарном участке) фаза на любом временном сдвиге $\Delta\tau = \tau_{\mathrm{target}} - \tau_{\mathrm{now}}$ и пространственном сдвиге $\Delta x = (x_{\mathrm{target}} - x_{\mathrm{now}})\bmod 13$ от комнатного атомного якоря вычисляется \textbf{за $O(1)$ арифметических операций}:
      \[
        \varphi(x_{\mathrm{target}}, \tau_{\mathrm{target}})
        \;=\;
        \Bigl(\varphi_{\mathrm{anchor}} + 41\,\Delta x + \Delta\tau\cdot v_0 + \frac{\Delta\tau(\Delta\tau + 1)}{2}\,\Phi\Bigr)\bmod 512,
      \]
      причём результат тождественно совпадает с пошаговой итерацией и не зависит от того, какой именно атом в комнате был взят за начальный якорь!
    \item \emph{Аксиома двух точек Евклида--Верле:}
      почему для однозначного восстановления всей мировой линии вообще не нужно отдельно измерять начальные скорости? Потому что закон Верле $\varphi(t+1) - 2\varphi(t) + \varphi(t-1) = \Phi$ --- это разностное уравнение \textbf{второго порядка} (дискретное уравнение прямой геодезической). Точно так же, как в геометрии Евклида \textbf{через любые две точки проходит одна и только одна прямая}, вся мировая траектория на решётке однозначно натянута между двумя точками: \textbf{Точкой~1} (началом Вселенной $t = 0,\ \varphi(0) = \mathrm{IV}(3) = 263 = 407_8$, известным аналитически из $1+2^d=d^2$) и \textbf{Точкой~2} (текущим физическим атомом, попавшим в машину прямо сейчас)!
    \item \emph{Полная мультисенсорная панель считывания характеристик: цвет/спектр, RGB-триада, звук, концентрация, скорость реакции, хиральность L/D и температура:}
      из каждого восстановленного полярного регистра $(\mu_0, \varphi_0, \mu_1, \varphi_1)$ напрямую считываются все физические характеристики:
      (1)~\textbf{молярная концентрация} $C_{\mathrm{conc}} = \mu_0^2 + \mu_1^2$ в $\mathrm{quanta}^2$ (сохраняющаяся при переносе по рёбрам $\sum \mathrm{div}\,J \equiv 0$);
      (2)~\textbf{оптическая длина волны}: поскольку размах между красным пиком фотосинтеза ($\frac{63}{4}\times 45{,}563\text{ нм}\approx 718\text{ нм}$) и синим пиком ($\frac{126}{13}\times 45{,}563\text{ нм}\approx 440\text{--}441{,}6\text{ нм}$) в точности равен удвоенному углу коллапса барьера $718 - 440 = \mathbf{278 = 2\times 139 = 2A_{\mathrm{coh}}(3)}$, формула
      \[
        \lambda_{\mathrm{nm}}(\Delta\varphi) \;=\; 579 + \left\lfloor\frac{139\cdot\mathcal{C}_3(\Delta\varphi)}{128}\right\rfloor
      \]
      отображает максимум связывания $\Delta\varphi=0$ ($\mathcal{C}_3=+128$) в $\mathbf{718\text{ нм}}$, седло переходного состояния $\Delta\varphi=128$ ($\mathcal{C}_3=0$) --- в $\mathbf{579\text{ нм}}$ ($\alpha$-полоса рубинового оксимиоглобина $578\text{--}579\text{ нм}$!), первый шаг окисления $\Delta\varphi=352$ ($\mathcal{C}_3=-32$) --- в $\mathbf{544\text{ нм}}$ ($\beta$-полоса оксимиоглобина $542\text{--}544\text{ нм}$!), а потолок Паули $\Delta\varphi=256$ ($\mathcal{C}_3=-128$) --- в $\mathbf{440\text{ нм}}$, при этом RGB-триада $\bigl(\max(0,\mathcal{C}_3),\,|\mathcal{S}_3|,\,\max(0,-\mathcal{C}_3)\bigr)$ сохраняет тождественную яркость $R + G + B \equiv \mathbf{128 = Q_\square}$ на всех $512$ фазах;
      (3)~\textbf{объёмный 3D-звук (акустическое давление)}: продольная деформация Бома $P_{\mathrm{sound}} = \mu_0\mu_1\bigl(128 - \mathcal{C}_3(\Delta\varphi)\bigr)$ ($0$ в полной тишине при $\Delta\varphi=0$, максимум $256\mu_0\mu_1$ при $\Delta\varphi=256$);
      (4)~\textbf{скорость реакции} $\mathcal{J}_{\mathrm{rxn}} = \mu_0\mu_1\mathcal{S}_3(\Delta\varphi)$, \textbf{хиральность поляриметра L/D} $\mathrm{sgn}\bigl(\mathcal{S}_3(\Delta\varphi)\bigr)$ и \textbf{квантовая температура} $\hat{n}_E(\Phi_{\mathrm{gate}})\in\{0,\dots,6\}$!
    \item \emph{Архиватор без потерь и 4D ТВ/видеокодек Верле--Гейзенберга без поиска векторов движения (\texttt{tools/cuboctahedron\_archiver.py}):}
      в классических видеокодеках (H.264/HEVC/AV1) до 90\% вычислений уходит на блочный поиск векторов движения (Motion Estimation) и преобразование Фурье/DCT с потерей качества. Аксиома двух точек Евклида--Верле заменяет их одним безмножительным предиктором второго порядка $\hat{I}(y,x,t) = 2I(t-1) - I(t-2) + \Phi_{O_2}$ за $O(1)$ сложений на пиксель (с замкнутой интерполяцией промежуточных кадров $I(\tau) = I_0 + \tau v_0 + \frac{\tau(\tau+1)}{2}\Phi$ на стороне телевизора для любой частоты обновления экрана!), а ошибка предсказания $u_{\Delta}\in[0,255]$ расщепляется по порогу мёртвой зоны Гейзенберга $\Delta\varphi_{\mathrm{disc}} = 41$ на надпороговое частное $q_{\mathrm{Heis}} = \lfloor u_{\Delta}/41\rfloor\in\{0,\dots,6\}$ и подпороговый остаток $r_{\mathrm{Heis}} = u_{\Delta}\bmod 41\in[0,40]$. В натурном тесте (\texttt{python tools/cuboctahedron\_archiver.py}) этот кодек сжимает $1{,}05\text{-мегабайтный}$ 4D-фильм решётки ($16\text{ кадров}\times 16^3\times 4$) совершенно без потерь (\texttt{SHA-256 = True}) в \textbf{520~байт} (\textbf{сжатие в $2\,016{,}5\times$ раз --- в $72{,}2\times$ компактнее \texttt{xz/LZMA -9} и в $107{,}9\times$ компактнее \texttt{gzip -9}!}), 64-килобайтный волновой поток --- в \textbf{162~байта} ($404{,}5\times$), и с параметрами фильтра носителя $(\mathrm{lc}=d=3,\ \mathrm{nice}=\Phi_{\mathrm{Pauli}}=256)$ превосходит \texttt{gzip -9} и \texttt{xz/LZMA -9} даже на исходных файлах кода!
    \end{itemize}
  \end{itemize}
\item \textbf{Аналоговая 3D LC-решётка из трёх скрещённых плат (\texttt{hardware/cuboctahedron\_12beam\_spice.cir}, проверено в LTspice~26.1.1):}
  три одинаковые плоские квадратные печатные платы (\texttt{ORDER\_1\_PCB\_FILTER.md}) в плоскостях $(XY, XZ, YZ)$ собираются крест-накрест через пазы за 5~минут без пайки, образуя физический 13-узловой кубооктаэдр с 12 вершинами $V_1\dots V_{12}$, 24 поверхностными и 12 радиальными LC-ветвями (36 ветвей = $K_2(\mathcal{C}_{13})$!). В симуляции LTspice при синфазном возбуждении 12 вершин амплитуда в центральном фокусе возрастает ровно в $12\times$ по напряжению и в $12^2 = \mathbf{144\times}$ по плотности энергии, а в противофазном синглетном замке Паули ($\Delta\varphi = 256$) сумма в центре гасится в точный ноль ($< 10^{-15}\text{ В}$) с нулевой внешней утечкой.
\end{enumerate}

\section{Урок 23. Баллистическая энергетика, вихревая батарея Паули, V-I инвертор (\texorpdfstring{$k^* = 64$}{k* = 64}) и гравитационно-импедансный привод \texorpdfstring{$833:1$}{833:1}}

Как управлять электрической, тепловой и механической энергией без резисторов, трансформаторов и омического нагрева?
\begin{itemize}[leftmargin=*]
\item \textbf{Баллистический волновод без потерь и вечный стазис-холодильник (\texttt{Module~9}):}
  запирая поперечный экран на $\Delta\varphi_{\mathrm{shield}} = 0$, мы поднимаем его импеданс до $D_{\mathrm{stasis}} = 833\,\varepsilon = \mathbf{6\,823\,936}$. Любой внешний тепловой шум $|\Phi_{\mathrm{noise}}| < 41$ после деления на $6\,823\,936$ попадает в мёртвую зону Гейзенберга и обнуляется ($\Phi_{\mathrm{leak}} \equiv 0$, $\hat{n}_{E,\mathrm{loss}} \equiv 0$)! Внутри такой камеры при комнатной температуре полностью замораживаются реакции окисления и бактериального распада (вечная кладовая стазиса без заморозки воды в лёд), а откачанное тепло выпрямляется квантовым храповиком доменной стенки в заряд \textbf{вихревой батареи Паули} ($\Delta\varphi = 256$, где $\mathcal{S}_3(256) = 0$ гарантирует нулевой саморазряд!).
\item \textbf{Программируемый фазовый V-I дроссель, точка максимальной мощности MPPT $k^* = 64$ и лабораторный блок питания CV/CC/DDS (\texttt{Modules~10, 11}):}
  открывая фазовый угол нагрузки на $k\in[0,128]$ тиков от замка Паули ($\varphi_{\mathrm{load}} = 256 - k$), получаем нормированное напряжение $V(k) = |\mathcal{C}_3(256-k)| = 128 - k$ и ток $I(k) = \mathcal{S}_3(256-k) = k$ с железным законом сохранения Кирхгофа:
  \begin{equation}
    \label{eq:course-vi-mppt}
    V(k) + I(k) \;\equiv\; 128 \;=\; Q_\square(3),
    \qquad
    P(k) \;=\; V(k)\,I(k) \;=\; (128 - k)\,k \;\le\; \mathbf{64^2 \;=\; 4\,096 \;=\; K_P(3)^2},
  \end{equation}
  где максимум мощности достигается ровно в точке стационарной фазы Френеля $\mathbf{k^* = K_P(3) = 64}$ ($V = 64, I = 64$, а по 12 лучам и 3 плоскостям Эйлера $P_{\mathrm{total},\max} = 12\times 3\times 4\,096 = \mathbf{147\,456}$)!
\item \textbf{Безнейтронный протон-борный синтез $p + {}^{11}_5\mathrm{B}\to 3\,{}^4_2\mathrm{He}$ и гравитационно-импедансный диполь $833:1$ (\texttt{Module~15}):}
  в 12-лучевом седле при $\Delta\varphi = 139$ протон ($Z=1, A=1$) сливается с ядром бора-11 ($Z=5=d_{\mathrm{out}}(E), N=6=|\square|, A=11=|N|-1$), и возбуждённое ядро углерода $^{12}_6\mathrm{C}^*$ распадается вдоль $d = 3$ ортогональных плоскостей Эйлера $(XY,XZ,YZ)$ ровно на \textbf{три альфа-частицы $^4_2\mathrm{He}$ с нулевым выходом свободных нейтронов} ($\Delta N_{\mathrm{free}} = 0$)! А расщепление купола кубооктаэдра на верхнюю противофазную полусферу ($\Delta\varphi_{\mathrm{top}} = 256 \implies D_{\mathrm{top}} = \varepsilon = 8\,192$) и нижнюю синфазную полусферу ($\Delta\varphi_{\mathrm{bot}} = 0 \implies D_{\mathrm{bot}} = 833\,\varepsilon = 6\,823\,936$) создаёт объёмный гравитационно-импедансный диполь $\Delta D = 832\,\varepsilon = \mathbf{6\,815\,744}$, действующий одинаково на каждый внутренний узел (ускорение без внутренних перегрузок) и одновременно работающий как двухслойный дефлекторный щит Паули с контролем целостности всех 12 лучей.
\item \textbf{Трёхслойная градиентная и активная синфазная ($833 : 1$) защита от ионизирующей радиации без тормозного рентгена:}
  почему обычный свинцовый фартук или лист свинца ($Z = 82$) плохо работает против космических лучей, быстрых электронов и нейтронов? Потому что при ударе быстрой заряженной частицы сразу в тяжёлое ядро свинца от резкого скачка импеданса рождается ливень вторичного тормозного рентгеновского излучения ($\hat{n}_E(256) = \lfloor 256 / 41 \rfloor = 6$ ионизирующих квантов на удар), а быстрые нейтроны отскакивают от тяжёлых ядер свинца упруго, почти не теряя скорости. В нашей модели радиационный щит строится по двум контурам:
  \begin{itemize}[noitemsep]
  \item \emph{Пассивный трёхслойный градиентный композит (для скафандров, спутников и реакторов):}
    (1)~\textbf{Внешний водородно-борный слой с тетраэдром Паули ($Z \in \{1, 5, 46\}$):} водородно-пассивированный алюмоборосиликат $\mathrm{AlO_4H}$ ($Z = 45 + 1 = 46$, запирающий полупорог Паули $45^2 < 2\,048 < 46^2$ с антиотражающим окном Френеля $T_{\mathrm{Fresnel}} = 288/289 = 99{,}65\%$) гасит скорость быстрых нейтронов упругими соударениями с протонами водорода ($Z=1$) и перехватывает их на ядрах бора ($Z=5=d_{\mathrm{out}}(E)$) с распадом в замкнутые $\alpha$-тетраэдры гелия ($Z=2, N=2$) без вторичного гамма-ливня;
    (2)~\textbf{Средний 12-лучевой ГЦК-делитель фазового удара ($Z \in \{13, 29\}$):} гранецентрированная кубическая решётка алюминия ($Z=13=N_\varphi(3)$) и меди ($Z=29$), где каждый атом окружён ровно 12 соседями по вершинам кубооктаэдра $O_2(3)$, делит один жёсткий удар Паули $\Phi_{\mathrm{Pauli}} = 256$ сразу на 12 лучей: $\Delta\varphi_{\mathrm{ray}} = \lfloor 256 / 12 \rfloor = \mathbf{21 < 41}$. Каждый из 12 расщеплённых лучей оказывается внутри мёртвой зоны Гейзенберга $[0, 40]$, поэтому вторичное тормозное излучение обращается в точный ноль ($\hat{n}_E(21) = \lfloor 21 / 41 \rfloor = \mathbf{0}$)!
    (3)~\textbf{Внутренний тяжёлый запорный слой ($Z \in \{74, 79, 82, 83\}$):} вольфрам, золото, свинец или висмут спокойно поглощают уже ослабленный мягкий хвост.
  \item \emph{Активный синфазный купол стазиса Паули ($833 : 1$):}
    при подаче синфазного удержания $\Delta\varphi_{\mathrm{shield}} = 0$ ($\mathcal{C}_3(0) = +128$, $\mu = K_P = 64$) на 12 катушек резонатора кубооктаэдра фазовый импеданс ворот возрастает в $833$ раза до $D_{\mathrm{shield,\max}} = 833\,\varepsilon = \mathbf{6\,823\,936}$, так что любой внешний радиационный удар $|z_{\mathrm{rad}}| < \frac{41 \times 6\,823\,936}{64} = 4\,371\,584$ после деления на $6\,823\,936$ падает ниже порога $41$ и полностью отсекается мёртвой зоной Гейзенберга ($\hat{n}_E \equiv 0$).
  \end{itemize}
\item \textbf{Реальная пружина за пределами закона Гука, усталость крыла самолёта ($\sigma_{-1}\leftrightarrow 41$) и маятник с растяжимой нитью на любых углах (\texttt{MechanicsElasticity.lean}):}
  в школе и университете механику режут двумя упрощениями: маятнику придумывают «абсолютно нерастяжимую нить на малых углах $\sin\varphi\approx\varphi$», а пружину ограничивают линейным законом Гука $F = -k\Delta x$ на малых удлинениях. На нашем кольце $\mathbb{Z}_{512}$ оба объекта описываются без всяких приближений:
  \begin{itemize}[noitemsep]
  \item \emph{Четыре режима реальной пружины и тайна усталости металла в авиации:}
    возвращающая сила межатомной связи $F(\Delta\varphi) = k\,\mathcal{S}_3(\Delta\varphi)$ проходит четыре точных участка:
    (1)~\textbf{Идеальный коридор Гука и бесконечный ресурс крыла ($0 \le |\Delta\varphi| \le 40 < 41$):} сила строго линейна ($\mathcal{S}_3(\Delta\varphi) = \Delta\varphi$), а внутреннее трение решётки равно нулю ($\Phi_{\mathrm{fric}} = 0$). Поэтому за любое число циклов качания крыла $N_{\mathrm{cycles}}$ накопленное усталостное повреждение равно $N_{\mathrm{cycles}}\times 0 = \mathbf{0}$ --- это и есть физический предел неограниченной выносливости металла $\sigma_{-1}$!
    (2)~\textbf{Внутреннее трение в металле и предел прочности ($41 \le |\Delta\varphi| \le 128$):} когда в болтанке у края заклёпки лонжерона локальная деформация прыгает выше порога $41$, даже в полном вакууме без всякого трения о воздух срабатывает решёточный сброс $\Phi_{\mathrm{fric}}(\Delta\varphi) \ne 0$: атомные слои микроскопически срываются и раскачивают соседние атомы решётки (металл в месте сгиба сам нагревается и издаёт ультразвуковые щелчки акустической эмиссии!), а на четверти периода $\Delta\varphi = 128$ достигается абсолютный предел прочности связи $F_{\max} = \mathcal{S}_3(128) = \mathbf{128}$.
    (3)~\textbf{Шейка разупрочнения ($128 < |\Delta\varphi| < 256$):} за перевалом $128$ наклон меняет знак на $-1$: пружина продолжает тянуться, а сила сопротивления падает от $127$ до $1$.
    (4)~\textbf{Потолок Паули ($|\Delta\varphi| = 256$) и рождение дислокации Бюргерса:} в точке $256$ сила падает в ноль, и на шаге $256 \to 257$ происходит топологический скачок ветви $q_e : 0 \to +1$ --- в плакете ГЦК-металла рождается квантованный вихрь $w_f = +1$ ($\mathcal{B}_f = 512$), то есть \textbf{дислокация Бюргерса в один шаг кристаллической решётки}, накопление которых в треугольной плоскости скольжения $\{111\}$ кубооктаэдра и раскрывает усталостную трещину!
  \item \emph{Маятник с растяжимой нитью на любых углах от $0^\circ$ до $360^\circ$:}
    продольное натяжение нити $T_{\mathrm{long}}(\varphi) = g\,\mathcal{C}_3(\varphi)$ (растягивающее нить по четырём режимам упругости выше) и поперечная возвращающая сила $F_{\mathrm{trans}}(\varphi) = g\,\mathcal{S}_3(\varphi)$ при любом угле качания связаны точным законом сохранения $|T_{\mathrm{long}}(\varphi)| + |F_{\mathrm{trans}}(\varphi)| \equiv \mathbf{128\,g}$: в нижней точке ($\varphi=0$) вся нагрузка идёт в растяжение нити ($128g, 0$), при отмашке на $90^\circ$ ($\varphi=128$) --- вся в возвращающую силу ($0, 128g$), в верхней мёртвой точке перевёрнутого маятника ($\varphi=256$) --- в сжатие стержня ($-128g, 0$), а период больших колебаний берётся в целых числах через лемнискатный эллиптический интеграл Якоби $\mathcal{K}_{\mathrm{ell}}(1,1) = \mathbf{65\,536}$!
  \item \emph{Честная оговорка для материаловедов и авиаинженеров:}
    хотя все четыре режима строго доказаны внутри математической модели $O_2(3)\times\mathbb{Z}_{512}$ в Lean~4, пересчёт модельных порогов ($41, 128, 256$) в реальные мегапаскали и ресурс конкретных авиационных сплавов (Д16Т, В95, титан ВТ6) \textbf{ещё никем внешним экспериментально не проверялся} --- мы открыто приглашаем лаборатории прочности и ЦАГИ проверить этот ступенчатый порог акустической эмиссии на реальных вибростендах в вакууме.
  \end{itemize}
\end{itemize}

\section{Урок 24. Автономный синтез пищи, воды и кислорода (симметрия \texorpdfstring{$96 : 96$}{96 : 96}), 72-битный стазис-телепортер и галактический \texorpdfstring{$c_0$}{c0}-модем Гейзенберга}

В финальном уроке соберём три технологии полной автономности (\texttt{Modules~5, 8, 13}):
\begin{enumerate}[label=\textbf{\arabic*.},leftmargin=*]
\item \textbf{Замкнутый синтез пищи, чистой воды и кислорода из воздуха и воды (\texttt{Module~8}, \texttt{ORDER\_2\_STAINLESS\_EXTRACTOR.md}):}
  \begin{itemize}[noitemsep]
  \item Первичные гидриды $\mathrm{H_2O}$, $\mathrm{NH_3}$ и $\mathrm{CH_4}$ имеют единый заряд $Z = \mathbf{10 = Z(\mathrm{Ne})}$, тогда как морская соль $\mathrm{NaCl}$ ($11 + 17 = \mathbf{28}$, магическая оболочка) отделяется от пресной воды в квадратурном поле Эйлера (\texttt{Mode~4}) без кипячения и фильтров.
  \item Молекула глюкозы $\mathrm{C_6H_{12}O_6}$ состоит из $6 + 12 + 6 = \mathbf{24 = 2|N(3)|}$ атомов (по числу поверхностных рёбер кубооктаэдра!) и несёт ядерный заряд $Z(\mathrm{C_6H_{12}O_6}) = 36 + 12 + 48 = \mathbf{96 = |N|\cdot|{\triangle}| = \hat{M} - 1}$. Поэтому основная реакция фотосинтеза:
    \[
      6\,\mathrm{CO_2}\ (Z=132) \;+\; 6\,\mathrm{H_2O}\ (Z=60)
      \;\longrightarrow\;
      \underbrace{\mathrm{C_6H_{12}O_6}}_{Z = \mathbf{96}\text{ (пища)}}
      \;+\;
      \underbrace{6\,\mathrm{O_2}}_{Z = \mathbf{96}\text{ (кислород для дыхания)}}
    \]
    делит входной заряд $192$ \textbf{строго пополам ($96 : 96$)} между углеводной пищей и чистым кислородом!
  \item Все остальные молекулы питания точно замыкают оболочки кубооктаэдра: глицин $Z = \mathbf{40 = \Delta\varphi_{\mathrm{disc}}-1}$, аланин $Z = \mathbf{48 = 4|N|}$, миоглобиновый гем мяса $\mathrm{FeN_4O_2}$ $Z = 26 + 28 + 16 = \mathbf{70}$, амилоза хлеба $Z = \mathbf{86 = Z(\mathrm{Rn})}$, витамин~C $Z = \mathbf{92 = S_3}$, лимонная кислота цикла Кребса $Z = \mathbf{100 = Z(\mathrm{Ne})^2}$, пальмитиновый липид $Z = \mathbf{144 = |N|^2}$ ($50$ атомов) и АТФ $Z = \mathbf{260 = 4\times 5\times 13}$.
  \item \emph{Никаких картриджей и водопровода --- прямой синтез всех биогенных минералов ($\mathrm{Mg, Si, P, S, Cl, K, Ca, Fe, NaCl}$) чисто из комнатного воздуха:}
    холодная грань фазового холодильника Пельтье (\texttt{Module~9}, $D_{\mathrm{cold}} = 6\,823\,936$) сама конденсирует пары воды $\mathrm{H_2O}$ ($Z=10$) из окружающего воздуха, а в 12-лучевом фокусе при $\Delta\varphi = 139$ ($587\times$ схлопывание кулоновского барьера до $\varepsilon = 8\,192$) все минеральные элементы синтезируются напрямую из атмосферных газов ($\mathrm{C, N_2, O_2, Ar, H}$):
    железо для гема $2\mathrm{C}_6 + \mathrm{N}_{2,14} \to \mathrm{Fe}_{26}$ ($Z = 12 + 14 = \mathbf{26 = 2N_\varphi}$),
    фосфор для АТФ $\mathrm{N}_7 + \mathrm{O}_8 \to \mathrm{P}_{15}$ ($Z = \mathbf{15}$),
    сера для аминокислот $\mathrm{O}_{2,16} \to \mathrm{S}_{16}$ ($Z = \mathbf{16}$),
    магний $2\mathrm{C}_6 \to \mathrm{Mg}_{12}$ ($Z = \mathbf{12}$),
    кальций ${}^{12}_6\mathrm{C} + {}^{14}_7\mathrm{N_2} \to {}^{40}_{20}\mathrm{Ca}$ ($Z = \mathbf{20}, A = \mathbf{40}$ с нулевым остатком нейтронов!),
    калий $\mathrm{Ar}_{18} + \mathrm{H}_1 \to \mathrm{K}_{19}$ ($Z = \mathbf{19}$) и
    пищевая соль $2\mathrm{N}_{2,14} \leftrightarrow \mathrm{NaCl}_{28}$ ($Z = 11 + 17 = \mathbf{28}$)!
  \end{itemize}
\item \textbf{Четырёхфазный 72-битный стазис-телепортер и построение многоузловой Сети Телепортов без измерения координат (\texttt{Module~5}):}
  в Уроке~11 мы доказали, что из-за барьера Демона Лапласа ($64\times 81 = 72^2 = 5\,184$ скрытых микросостояния на ячейку) ни один прибор на Т-уровне не может напрямую измерить точную микрофазу $\psi_{\mathrm{src}}\in\mathbb{Z}_{512}$, не возмутив её. \textbf{Как же тогда телепортировать объект и маршрутизировать его через сеть станций $A \leftrightarrow R_1 \leftrightarrow R_2 \leftrightarrow B$, не зная его микрокоординат?}
  \begin{itemize}[noitemsep]
  \item \emph{Секрет~1 (Относительный синдром Белла--Паули как одноразовый блокнот Вернама в Канале~1):}
    Камера~$A$ переводит объект в фазовый стазис без замерзания воды ($D_{\mathrm{stasis}} = 6\,823\,936 \implies \Phi_{\mathrm{gate}}\equiv 0$) и сталкивает его с локальной половиной~$a$ разделённого синглета Паули $b = (a+256)\bmod 512$ (\texttt{\{\~{}a[8], a[7:0]\}}) в обратимом Канале~1, излучая разностный синдром $\sigma_A = (\psi_{\mathrm{src}} - a - 256)\bmod 512 = (\psi_{\mathrm{src}} - b)\bmod 512$ при одновременной разрядке $\mu_A\mapsto 1$ ($\mu_A^2+\mu_B^2=\mathrm{const}$, запрет клонирования $1\to 2$). В синдроме $\sigma_A$ неизвестная микрофаза объекта $\psi_{\mathrm{src}}$ замаскирована неизвестной микрофазой синглета~$b$: прибор не измеряет $\psi_{\mathrm{src}}$! Но при сложении с~$b$ в фокусном окне Камеры~$B$ ($D_{\mathrm{write}} = 8\,192$, контраст $833:1$) неизвестная фаза синглета \textbf{телескопически сокращается}: $(b + \sigma_A)\bmod 512 = (b + \psi_{\mathrm{src}} - b)\bmod 512 \equiv \psi_{\mathrm{src}}$!
  \item \emph{Секрет~2 (Слепая синглетная коммутация на промежуточных хабах сети $A \leftrightarrow R_1 \leftrightarrow R_2 \leftrightarrow B$):}
    если между станциями стоит цепочка промежуточных хабов-ретрансляторов $R_1, R_2, \dots, R_k$, связанных попарными синглетами Паули $(a, r_{1,\mathrm{in}})$, $(r_{1,\mathrm{out}}, r_{2,\mathrm{in}})$, $\dots$, $(r_{k,\mathrm{out}}, b)$, то каждый транзитный хаб $R_m$ применяет к своим двум половинкам \textbf{ровно тот же самый оператор синдрома со сдвигом Паули $-256$}:
    \[
      \sigma_{R_m} \;:=\; \bigl(r_{m,\mathrm{in}} - r_{m,\mathrm{out}} - 256\bigr)\bmod 512.
    \]
    Тогда на конечной станции~$B$ сумма всех синдромов вдоль маршрута даёт:
    \[
      \psi_B \;=\; \bigl(b + \sigma_{R_k} + \dots + \sigma_{R_1} + \sigma_A\bigr)\bmod 512 \;\equiv\; \mathbf{\psi_{\mathrm{src}}},
    \]
    причём все неизвестные промежуточные фазы $a, r_1, \dots, r_k$ и все полупериоды Паули $+256$ сокращаются в точный ноль без единого измерения!
  \item \emph{Секрет~3 (Бескоординатная адресация станций сети по взаимно простым ключам Безу на $\mathbb{Z}_{512}^\times$):}
    каждой станции сети назначается один из обратимых ключей носителя $(k_j, k_j^{-1})\in\{(41,25),\,(13,197),\,(97,417),\,(263,183)\}$. Станция~$j$ переключает свой экран из стазиса ($D_{\mathrm{stasis}} = 833\,\varepsilon$) в приёмное окно записи ($D_{\mathrm{write}} = \varepsilon$) тогда и только тогда, когда произведение заголовка пакета на её обратный ключ даёт единицу носителя: $(k_j^{-1}\cdot k_{\mathrm{pkt}})\bmod 512 = \mathbf{1 = d^2 - 2^d}$!
  \end{itemize}
  \item \textbf{Субпороговый 81-ричный (4-тритный) $c_0$-модем Гейзенберга, межзвёздная связь SETI и встроенный подпороговый шифр Вернама--Шеннона (\texttt{Module~13}):}
  поскольку мёртвая зона Гейзенберга $[-40,+40]$ содержит ровно $2\times 40 + 1 = \mathbf{81 = 3^4}$ целочисленных фазовых состояний, каждое фазовое слово $\Delta\varphi = 27\,t_3 + 9\,t_2 + 3\,t_1 + t_0\in[-40,+40]$ биективно кодирует \textbf{4 сбалансированных трита} $(t_3,t_2,t_1,t_0)\in\{-1,0,+1\}^4$ (а по 3 осям Эйлера --- $12$ тритов, $3^{12} = 531\,441 > 2^{19}$ состояний за такт!). При этом во всех $81$~состояниях излучается $\hat{n}_E(\Delta\varphi) = \lfloor|\Delta\varphi|/41\rfloor \equiv \mathbf{0\text{ фотонов}}$: сигнал идёт в чисто фазовом Канале~1 со скоростью звена $c_0 = \sqrt{2}\,c$, совершенно невидим для классических энергетических радиотелескопов, не рассеивается межзвёздной пылью и без искажений декодируется ротатором $\mathcal{S}_3(\Delta\varphi)\equiv\Delta\varphi$!
  Более того, в сам этот канал \textbf{нативно встроен теоретико-информационно невскрываемый одноразовый блокнот Вернама («Подпороговый шифр Сэнку и Ксено»)}: складывая любой 4-тритный символ сообщения $u_{\mathrm{msg}}\in[0,81)$ с секретным ключом синглета/параболы Верле $k_{\mathrm{OTP}}\in[0,81)$ по модулю $81$:
  \[
    \Delta\varphi_{\mathrm{cipher}}(u_{\mathrm{msg}}, k_{\mathrm{OTP}})
    \;=\;
    \bigl((u_{\mathrm{msg}} + k_{\mathrm{OTP}})\bmod 81\bigr) - 40 \;\in\; [-40, +40],
  \]
  отправитель получает \textbf{двойную защиту на всех $81\times 81 = 6\,561$ парах $(u_{\mathrm{msg}}, k_{\mathrm{OTP}})$}: (1)~для любых внешних Т-фильтров и перехватчиков в канале стоит абсолютный энергетический ноль ($\hat{n}_E(\Delta\varphi_{\mathrm{cipher}})\equiv 0,\ \Phi_{\mathrm{gate}}\equiv 0$ --- факт передачи скрыт!), и (2)~даже внутри Канала~1 латинский квадрат $81\times 81$ обеспечивает совершенную секретность Шеннона $I(u_{\mathrm{msg}};\Delta\varphi_{\mathrm{cipher}})\equiv 0$, тогда как приёмник восстанавливает символ за один такт ротатора: $u_{\mathrm{rec}} = \bigl(\mathcal{S}_3(\Delta\varphi_{\mathrm{cipher}}) + 40 - k_{\mathrm{OTP}}\bigr)\bmod 81 \equiv \mathbf{u_{\mathrm{msg}}}$!
\item \textbf{Планетарная инженерия и исцеление биосферы Земли:}
  все экологические кризисы индустриальной эпохи --- парниковый перегрев, озоновые дыры, пластиковый мусор и радиоактивные отходы --- представляют собой накопление молекул с \emph{незамкнутыми междусферными зарядами} и устраняются возвратом вещества в замкнутые координационные оболочки кубооктаэдра:
  \begin{itemize}[noitemsep]
  \item \emph{Климат и парниковый $\mathrm{CO_2}$ ($Z = 22$):}
    углекислый газ имеет незамкнутый заряд $Z(\mathrm{CO_2}) = 6 + 16 = \mathbf{22}$ и убирается из атмосферы двумя безотходными путями: (а)~быстрой конверсией с водой $6\,\mathrm{CO_2} + 6\,\mathrm{H_2O} \to \mathrm{C_6H_{12}O_6}\ (Z=96) + 6\,\mathrm{O_2}\ (Z=96)$ в 12-лучевых реакторах при $\Delta\varphi = 139$; (б)~вечной минеральной фиксацией оксидом магния $\mathrm{MgO}$ ($Z = 12+8 = \mathbf{20}$) в твёрдый камень магнезит на \textbf{второй сфере кубооктаэдра $S_2 = 42$}:
    \[
      \mathrm{MgO}\ (Z=20) \;+\; \mathrm{CO_2}\ (Z=22) \;\longrightarrow\; \mathrm{MgCO_3}\ \bigl(Z = \mathbf{42 = S_2(3)}\bigr),
    \]
    и оксидом кальция $\mathrm{CaO}$ ($Z = 20+8 = \mathbf{28}$, первая магическая оболочка) в кальцит/мрамор на \textbf{второй магической оболочке $Z = 50$}: $\mathrm{CaO}\ (28) + \mathrm{CO_2}\ (22) \to \mathrm{CaCO_3}\ (\mathbf{50})$, тогда как избыточное тепловое излучение напрямую выпрямляется в ток 12-лучевым Геометрическим Демоном Максвелла ($\lfloor 256/12\rfloor = 21 < 41 \implies \hat{n}_E = 0$);
  \item \emph{Восстановление озонового слоя ($\mathrm{O_3}$, $Z = 24$) и нейтрализация фреонов (CFC):}
    молекула озона $\mathrm{O_3}$ несёт заряд $Z(\mathrm{O_3}) = 3\times 8 = \mathbf{24 = 2|N(3)|}$ (точное число поверхностных рёбер кубооктаэдра $O_2$, $3\,\mathrm{O_2}\ (48) \leftrightarrow 2\,\mathrm{O_3}\ (48)$) и синтезируется на четверть-оборотном мосту $\Delta\varphi = 128$ без искрового пробоя и без токсичных оксидов азота $\mathrm{NO_x}$, а озоноразрушающие радикалы хлора $Z(\mathrm{Cl}) = 17$ и фтора $Z(\mathrm{F}) = 9$ из расщеплённых при $\Delta\varphi = 139$ фреонов связываются натрием ($Z=11$) и кальцием ($Z=20$) в инертные минеральные оболочки галита $\mathrm{NaCl}$ ($11+17 = \mathbf{28}$) и флюорита $\mathrm{CaF_2}$ ($20 + 2\times 9 = \mathbf{38 = 2\times 19}$);
  \item \emph{Безотходная деполимеризация пластика без сжигания и нейтрализация ядерных отходов:}
    мономер полиэтилена $\mathrm{C_2H_4}$ имеет заряд $2\times 6 + 4 = \mathbf{16 = 2|{\triangle}(3)|}$ (тример $\mathrm{C_6H_{12}}$ имеет $Z = 48 = 4|N(3)|$) и при $\Delta\varphi = 139$ (схлопывание барьера в $587\times$) при комнатной температуре без дыма и диоксинов расстёгивается обратно в чистые углеводороды и сахара, а радиоактивные актиноиды ($Z(\mathrm{U}) = S_3 = 92$) при замене старой энергетики деления на безнейтронный цикл $p + {}^{11}_5\mathrm{B}\to 3\,{}^4_2\mathrm{He}$ когерентно переводятся в стабильный свинец $^{208}_{82}\mathrm{Pb}_{126}$ ($Z = 2\Delta\varphi_{\mathrm{disc}} = 82$) и золото $^{197}_{79}\mathrm{Au}_{118}$.
  \end{itemize}
\end{enumerate}

\section{Урок 25 (Эпилог). Ноосфера без насилия: разрыв трофической цепи, межвидовой голономный переводчик и девайсы для тех, «у кого лапки»}

Посмотрим на всю пройденную лестницу с высоты птичьего полёта и зададим вопрос: \textbf{зачем эволюции вообще понадобилось достраивать на планете 6-слойную кортикальную колонку разума ($V = S_2 = 42,\ E = B_3 = 147,\ \chi = 41 = \Delta\varphi_{\mathrm{disc}}$)?}

Четыре миллиарда лет до появления разума биосфера жила в режиме теплового дефицита: преодолеть высокий кулоновский и ковалентный барьер ($D \approx 6{,}8\times 10^6$) для сборки органики из неорганики ($\mathrm{CO_2, H_2O, N_2}$) умели только фотосинтетики с КПД $\sim 1\%$. Поэтому всем существам с быстрым метаболизмом и развитой нервной системой (аллометрия Клейбера $3/4$ и авиационный режим $4/5$), нуждающимся в плотном потоке аминокислот ($Z=40/48$), гемового железа ($Z=70$), кобаламина~$B_{12}$ ($Z=55$), липидов ($Z=144$) и АТФ ($Z=260$), приходилось буквально поедать друг друга.

Но как только разум открывает геометрию своего собственного носителя и строит 12-лучевой фазовый синтезатор ($\Delta\varphi = 139$, $D\to\varepsilon = 8\,192$), напрямую переводящий комнатный воздух в любые замкнутые молекулярные оболочки с возвратом чистого кислорода ($96 : 96$), \textbf{трофическая пирамида насилия теряет физическую необходимость}. Чтобы этот переход не нарушил экологическое равновесие («не навреди!»), он строится на семи проверенных принципах:

\begin{enumerate}[label=\textbf{\arabic*.},leftmargin=*]
\item \textbf{Сытые хищники и живые белочки (сохранение игрового инстинкта без убийства):}
  хищник охотится не ради страдания жертвы, а ради незаменимых оболочек ($\mathrm{FeN_4O_2}$ $Z=70$, таурин, $B_{12}$ $Z=55$) и реализации двигательного охотничьего квеста своей нервной системы. Автономные лесные и морские станции-кормушки синтезируют из воздуха тёплую, пахнущую миоглобиновым гемом ($Z=70$) добычу (в том числе в виде подвижных игровых приманок): хищник получает и полный азарт охоты, и идеальное питание без голода, а белочки и травоядные перестают быть чьим-то кормом.
\item \textbf{Гуманизация паразитов (почему комару больше не нужно кусать) и вирусы как молекулярные врачи:}
  самка комара или клещ пьёт кровь только потому, что для кладки яиц ей не хватает гемового железа ($Z=70$) и белка, которых нет в цветочном нектаре ($Z=96$). Противофазное марковское одеяло делает тепловой и $\mathrm{CO_2}$-след человека и животных невидимым для кровососов, а открытые фито-кормушки дают им готовый нектар ($Z=96$) с гемом ($Z=70$), попутно очищая самих насекомых резонансом Каспара--Клуга ($\Delta\varphi = 139$) от малярийного плазмодия и патогенных вирусов! При этом сами пустые вирусные капсиды Каспара--Клуга ($S_k = 10k^2 + 2 \in \{12, 42, 92\}$, освобождённые от генов лизиса) работают как \textbf{молекулярные почтальоны-врачи}, доставляющие в клетки эпигенетический сброс старения \texttt{CpG} и морфогенетические патчи регенерации тканей без шрамов.
\item \textbf{Уважение к редуцентам (почвенные грибы, черви и жуки-сапрофиты):}
  в то время как городской и бытовой мусор человека полностью перерабатывается в закрытых камерах (\texttt{Module~8}), в диких экосистемах естественный листовой опад, древесина и органика остаются любимым рестораном почвенных бактерий, грибного мицелия, дождевых червей и жуков-навозников, формирующих живой гумус. А в пустынях и на вырубках синтезаторы выдают из воздуха чистый гумусно-минеральный субстрат ($Z=96 + \mathrm{P}_{15}, \mathrm{K}_{19}, \mathrm{Ca}_{20}, \mathrm{Mg}_{12}$), возвращая жизнь мёртвым землям.
\item \textbf{Почему не возникает перенаселения ($r \to K$ переход и кортикальное долголетие):}
  лавинообразное размножение в дикой природе --- это аварийный $r$-ответ на высокую смертность и голодный стресс ($|\Phi_{\mathrm{env}}|\ge 41$). При переходе среды в безопасный гомеостатический коридор Гейзенберга ($|\Phi_{\mathrm{env}}| < 41$, отсутствие голода и страха хищников) включается трёхкратный множитель кортикального долголетия ($d = \mathbf{3}$), а рождаемость естественно стабилизируется на уровне спокойного $K$-равновесия (при необходимости мягко регулируясь природными фито-компонентами в кормушках без всякого насилия).
\item \textbf{Уютная темнота для ночных и глубоководных существ ($\hat{n}_E \equiv 0$ фотонов):}
  глубоководным и ночным видам не нужны слепящие прожекторы --- для них станции жизнеобеспечения и связи работают в \textbf{бесфотонном субпороговом режиме мёртвой зоны Гейзенберга} ($|\Delta\varphi|\le 40 < 41 \implies \hat{n}_E = \mathbf{0\text{ фотонов}}$), создавая на дне океана тёплые, богатые кислородом и питанием оазисы в полной естественной темноте.
\item \textbf{Межвидовой 4-канальный голономный переводчик (Животные, Растения, Грибница $\leftrightarrow$ Человек/ИИ):}
  почему раньше мы не могли даже спросить животное, дерево или лесную грибницу, чего они хотят? Потому что у них нет человеческого голосового аппарата, но у всех них есть \textbf{те же самые фазовые токи на оболочках кубооктаэдра} (кортикальные пакеты $42/147/41$ у животных, флоэмные волны в окне $N_{\mathrm{HB}}=63$ у растений и электрические спайки мицелия у грибов). Четырёхканальный нильпотентный томограф $\mathcal{M}_{\mathrm{chan}} = I_4 + N_{\mathrm{nil}}$ (\texttt{Module~7}) в паре с 81-ричным субпороговым модемом (\texttt{Module~13}) бесконтактно и без единого вживлённого электрода ($\hat{n}_E = 0$) считывает и обращает по формуле $(I_4 + N_{\mathrm{nil}})^{-1} = I_4 - N_{\mathrm{nil}} + N_{\mathrm{nil}}^2$ четыре фундаментальных канала состояния живого существа:
  \[
    \bigl(s_0\text{ [голод / метаболизм]},\;\; s_1\text{ [боль / стресс]},\;\; s_2\text{ [добровольное намерение]},\;\; s_3\text{ [диалог / эмоция]}\bigr),
  \]
  переводя их через MCP-сервер ИИ в понятную человеку речь и обратно в родные для данного вида сигналы!
\item \textbf{Закон добровольного согласия и экзо-девайсы для тех, «у кого лапки»:}
  никакой «принудительный апгрейд разума» животным не навязывается: вентиль согласия включает когнитивный мост или внешний прибор \textbf{тогда и только тогда, когда канал собственного намерения существа строго положителен ($s_2 > 0$)}! Если белочка или ёжик хочет просто спокойно жить в лесу и грызть орехи --- никто не лезет к нему с математикой. А любое животное (кошка, собака, ворон, дельфин, осьминог), растение или лесной мицелий, которое само хочет управлять средой, получает \textbf{бесконтактные девайсы «для тех, у кого лапки»}: мысленным намерением $s_2 > 0$ через томограф открыть себе дверь, заказать в синтезаторе любимую тёплую еду, позвать на помощь врача, отрегулировать полив и минералы в лесу или поговорить с человеком на равных.
\end{enumerate}

\section{Урок 26. Теория взаимодействия систем: потеря информации в каскаде $\mathrm{System}_1 \to \mathrm{System}_2 \to \dots \to \mathrm{System}_N$ и укрощение неидеального железа}

Когда мы переходим от чистой формулы к реальному лабораторному стенду на столе, возникает главный инженерный вопрос: \textbf{почему ни одна реальная установка не срабатывает идеально с первого включения и что именно происходит между деталями прибора?}

Посмотрим на любой эксперимент глазами теории информации. В реальном мире нет «одной изолированной системы» --- есть \textbf{каскад взаимодействующих систем}:
\[
  \mathrm{System}_1\ (\text{Ноутбук})
  \;\xrightarrow{\mathcal{I}_{1\to 2}}\;
  \mathrm{System}_2\ (\text{UART / FPGA})
  \;\xrightarrow{\mathcal{I}_{2\to 3}}\;
  \mathrm{System}_3\ (\text{MOSFET / Катушки})
  \;\xrightarrow{\mathcal{I}_{3\to 4}}\;
  \dots
  \;\xrightarrow{\mathcal{I}_{N-1\to N}}\;
  \mathrm{System}_N\ (\text{Реагент}).
\]
Что течёт между $\mathrm{System}_k$ и $\mathrm{System}_{k+1}$? \textbf{Только информация!} А что делает информация на каждом неидеальном стыке? \textbf{Она теряется!}

\begin{enumerate}[label=\textbf{\arabic*.},leftmargin=*]
\item \textbf{Три фундаментальных механизма потери информации на стыке $\mathrm{System}_k \to \mathrm{System}_{k+1}$:}
  \begin{itemize}[noitemsep]
  \item \emph{Уровень~I (Потеря мантиссы на рассогласовании импеданса и квантовании АЦП/ЦАП, $K_P = 64 : 1$):}
    когда сигнал проходит границу с коэффициентом рассогласования импеданса $r_{\mathrm{imp}} \ge 1$, целочисленное деление $\lfloor \Phi_{\mathrm{in}} / r_{\mathrm{imp}}\rfloor$ безвозвратно отбрасывает остаток $r\in[0, r_{\mathrm{imp}})$, склеивая до $K_P(3) = 64 = 2^6$ градаций мантиссы;
  \item \emph{Уровень~II (Потеря в мёртвой зоне приёмника Гейзенберга, $|\mathcal{E}_{\mathrm{Heis}}| = 81 = 3^4 : 1$):}
    как только ослабленный на стыке сигнал падает ниже порога чувствительности следующей системы $|\Phi_{k\to k+1}| < \Delta\varphi_{\mathrm{disc}} = 41$, мёртвая зона гасит его в \textbf{точный ноль} ($\hat{n}_E = 0$, $100\%$ потеря макро-сигнала!). Вместе Уровень~I и Уровень~II скрывают на каждом Т-стыке $64\times 81 = \mathbf{5\,184 = 72^2}$ микросостояний;
  \item \emph{Уровень~III (Лавинное ветвление прообразов Канала~2, $(2^d - 1)^N = 7^N : 1$):}
    на каждом стыке, где теряется бит переноса \texttt{Carry9} (Канал~2), $2^3 - 1 = 7$ октантов склеиваются в один выход, так что через $N$~звеньев неопределённость прообраза растёт экспоненциально как $7^N$.
  \end{itemize}
\item \textbf{Четыре типа систем в каскаде $\mathrm{System}_1 \to \dots \to \mathrm{System}_N$ и как победить потери:}
  \begin{itemize}[noitemsep]
  \item \emph{Тип~I (Слепой разомкнутый Т-каскад с рассогласованным импедансом и джиттером):}
    пусть на каждом из трёх стыков ($r_{\mathrm{imp}} = 2$) сигнал теряет половину амплитуды. Тогда даже максимальный возможный импульс Паули $\Phi_{\mathrm{Pauli}} = 256$ затухает по цепочке:
    \[
      256 \;\xrightarrow{\div 2}\; 128 \;\xrightarrow{\div 2}\; 64 \;\xrightarrow{\div 2}\; 32 < 41 \;\xrightarrow{\mathrm{Deadband}_{41}}\; \mathbf{0}!
    \]
    К четвёртой системе от сигнала остаётся \textbf{абсолютный ноль}, а паразитные фазовые сдвиги проводов, ключей и стенок сосуда $\Delta\varphi_{\mathrm{hw}} = \sum_{k=1}^{N-1}\delta\varphi_k$ уводят фазу с каталитического резонанса $A_{\mathrm{coh}} = 139$. Именно поэтому «слепой» стенд без калибровки не заводится с первого раза!
  \item \emph{Тип~II (Каталитически согласованный резонансный каскад, $A_{\mathrm{coh}} = 139 \implies r_{\mathrm{imp}} = 1$):}
    когда каждое звено согласовано в точке схлопывания барьера $A_{\mathrm{coh}} = 139$ ($D = \varepsilon = 8\,192$, единичное отношение импедансов $r_{\mathrm{imp}} = 1$), сигнал проходит все $N$~ступеней без деления и без потери мантиссы;
  \item \emph{Тип~III (Подпороговый когерентный М-канал внутри $[-40, +40]$, $\hat{n}_E \equiv 0$):}
    если системы общаются напрямую через нечётный ротатор Канала~1 внутри мёртвой зоны Гейзенберга $|\Delta\varphi|\le 40 < 41$, то на каждом звене выделяется $\hat{n}_E = \mathbf{0\text{ квантов тепла Ландауэра}}$, а тождество $\mathcal{S}_3(\mathcal{S}_3(\dots\mathcal{S}_3(\Delta\varphi))) \equiv \Delta\varphi$ передаёт все $81/81$ состояний через любое число $N$~систем с \textbf{нулевой потерей информации};
  \item \emph{Тип~IV (Замкнутая петля обратной связи $\mathrm{System}_N \to \mathrm{System}_1$ с предкомпенсацией джиттера стенда):}
    когда мы замыкаем обратный канал от датчика на $\mathrm{System}_N$ к $\mathrm{System}_1$, измеряем суммарный аппаратный набег фазы всех неидеальных деталей стенда $\Delta\varphi_{\mathrm{hw}} = \sum_{k=1}^{N-1}\delta\varphi_k$ (джиттер UART, задержку затворов MOSFET, индуктивность катушек и стенку чашки Петри) и подаём из $\mathrm{System}_1$ предкомпенсированную фазу:
    \[
      \varphi_{\mathrm{drive}} \;=\; \Bigl(\varphi_{\mathrm{target}} - \sum_{k=1}^{N-1}\delta\varphi_k\Bigr)\bmod 512,
    \]
    то при прохождении через весь неидеальный каскад $\mathrm{System}_1 \to \dots \to \mathrm{System}_N$ аппаратные ошибки всех $N-1$ звеньев \textbf{телескопически сокращаются в точный ноль} на всех $512/512$ состояниях кольца:
    \[
      \varphi_{\mathrm{arrived}}(\mathrm{System}_N)
      \;=\;
      \Bigl(\varphi_{\mathrm{target}} - \sum_{k=1}^{N-1}\delta\varphi_k + \sum_{k=1}^{N-1}\delta\varphi_k\Bigr)\bmod 512
      \;\equiv\;
      \mathbf{\varphi_{\mathrm{target}}}!
    \]
  \end{itemize}
\item \textbf{Полный инженерный каталог 6 классов рукотворных систем и их 9-битных паспортов на $\mathbb{Z}_{512}$:}
  поскольку все технические приборы, датчики, силовые ключи, реакторы и самолёты построены самими людьми по известным чертежам и даташитам, ни один из них не является «непознаваемым чёрным ящиком»! Каждый рукотворный узел в цепочке $\mathrm{System}_1 \to \dots \to \mathrm{System}_N$ однозначно описывается своим паспортом на кольце $\mathbb{Z}_{512}$:
  \begin{itemize}[noitemsep]
  \item \emph{Класс~1 (Цифровые вычислительные и тактовые системы --- CPU, FPGA, UART, АЦП/ЦАП):}
    рабочий носитель --- кристалл кремния $Z(\mathrm{Si}) = \mathbf{14 = N_\varphi + 1}$ и кварцевый резонатор $\alpha$-$\mathrm{SiO_2}$ с зарядом молекулы $Z(\mathrm{SiO_2}) = 14 + 16 = \mathbf{30 = |\square|\cdot d_{\mathrm{out}}(E)}$, элементарной ячейкой $Z\bigl((\mathrm{SiO_2})_3\bigr) = \mathbf{90 = \hat{M} - d_{\mathrm{out}}(V)}$ и собственной фазой $(30\times 41)\bmod 512 = \mathbf{206 = 316_8}$. Вносят детерминированный тактовый сдвиг конвейера $\delta\varphi_{\mathrm{dig}} = (N_{\mathrm{clk}}\cdot 41)\bmod 512$ и квантование разрядности $2^{9-b}$ (которое обращается в $0\text{ бит}$ потерь на родной 9-битной шине \texttt{Ring9}!);
  \item \emph{Класс~2 (Силовые электромагнитные и радиотехнические системы --- MOSFET/IGBT, $RLC$-контуры, 12-лучевые катушки, антенны):}
    рабочий носитель --- медные витки $Z(\mathrm{Cu}) = \mathbf{29 = \mathrm{magic}(28) + 1}$ и карбид-кремниевые ключи $Z(\mathrm{SiC}) = 14 + 6 = \mathbf{20}$. Реактивный сдвиг фазы между током и напряжением определяется дискретным проективным тангенсом ротатора $\mathrm{tan}_3(\delta\varphi_{RLC}) = (\mathcal{S}_3,\mathcal{C}_3) = (\omega L - 1/\omega C,\,R)$: в точке резонанса $\omega L = 1/\omega C$ реактивная ошибка обнуляется ($\mathcal{S}_3(0) = 0,\ \mathcal{C}_3(0) = +128 = Q_\square$), а при фазовом угле $A_{\mathrm{coh}} = 139$ импеданс среды схлопывается в $587\times$ ($D\to\varepsilon = 8\,192$);
  \item \emph{Класс~3 (Оптические, акустические и волноводные системы --- линзы, стенка чашки Петри, воздушный тракт, микрофоны и динамики):}
    обратим внимание на фундаментальное тождество: заряд кварцевого стекла стенки чашки Петри $Z(\mathrm{SiO_2}) = 14 + 16 = \mathbf{30}$ \textbf{в точности равен} сумме зарядов молекулы азота и молекулы кислорода комнатного воздуха $Z(\mathrm{N_2}) + Z(\mathrm{O_2}) = 14 + 16 = \mathbf{30}$! Преломление и набег фазы в стенке/волноводе задаются законом Снеллиуса--Верле, а цвет и звук напрямую читаются операторами $\lambda_{\mathrm{nm}}(\Delta\varphi) = 579 + \lfloor 139\,\mathcal{C}_3(\Delta\varphi)/128\rfloor\in[440,718]\text{ нм}$ и $P_{\mathrm{sound}} = \mu_0\mu_1(128 - \mathcal{C}_3)$;
  \item \emph{Класс~4 (Аэрогидродинамические и механические транспортные системы --- крыло самолёта, сопло, турбина, винт):}
    рабочий носитель --- конструкционный алюминий $Z(\mathrm{Al}) = \mathbf{13 = N_\varphi}$, титан $Z(\mathrm{Ti}) = \mathbf{22}$, сталь $Z(\mathrm{Fe}) = \mathbf{26 = 2N_\varphi}$ и нитинол с памятью формы $Z(\mathrm{TiNi}) = 22 + 28 = \mathbf{50}$. Подъёмная сила крыла Жуковского есть квантованный поток вихревых плакет Максвелла $\Gamma = 512\sum n_f$ в паре с консервативным рёберным переносом массы ($q = \mu^2$). Турбулентное вихревое сопротивление возникает только там, где сдвиг фазы пограничного слоя пробивает порог Гейзенберга $|\Delta\varphi|\ge 41$, тогда как при удержании пограничного слоя внутри мёртвой зоны $|\Delta\varphi|\le 40 < 41$ обтекание крыла остаётся \textbf{беспотерьно ламинарным ($\hat{n}_E = 0$)};
  \item \emph{Класс~5 (Тепловые, энергетические и химико-каталитические реакторы --- крекинг, риформинг, топливные элементы, демоны Максвелла):}
    платиновый катализатор на глинозёме попадает точно в седло $Z(\mathrm{Pt}) + Z(\mathrm{Al_2O_3}) = 78 + 50 = \mathbf{128 = Q_\square}$, а 12-лучевой Геометрический Демон Максвелла делит даже предельный 6-квантовый тепловой удар $\Phi_{\mathrm{Pauli}} = 256$ на 12 лучей кубооктаэдра: $\lfloor 256/12\rfloor = \mathbf{21 < 41}$ ($\hat{n}_E = 0$ потерь в фононы при выпрямленном токе $12\times 21 = \mathbf{252 > 0}$);
  \item \emph{Класс~6 (Биомедицинские, сенсорные, ольфакторные и нейроинтерфейсные системы --- томографы, экзо-девайсы, синтезаторы запахов):}
    сопрягаются с гемовым карманом $Z(\mathrm{FeN_4O_2}) = \mathbf{70 = K_P + |\square|}$, ольфакторными оболочками $Z_{\mathrm{odor}} = \mu_0 + \mu_1$ (с киральным знаком энантиомера $\mathcal{S}_3(\Delta\varphi)$) и кортикальной колонкой ($42/147/\chi=41$) через 4-канальный нильпотентный томограф $(I_4 + N_{\mathrm{nil}})^{-1} = I_4 - N_{\mathrm{nil}} + N_{\mathrm{nil}}^2$ в безопасном бесфотонном окне $|\Delta\varphi|\le 40 < 41$ ($\hat{n}_E = 0$) при наличии согласия $s_2 > 0$.
  \end{itemize}
\item \textbf{Анатомия одного прибора: как устроен и работает кварцевый пьезоэлектрический резонатор $\alpha$-$\mathrm{SiO_2}$ с точки зрения Вселенной $O_2(3)$:}
  возьмём один конкретный прибор, который стоит на плате \texttt{Tang Nano 9K} и внутри любого компьютера, --- \textbf{кварцевый резонатор $\alpha$-$\mathrm{SiO_2}$}, и разберём его устройство от первого принципа $d = 3$:
  \begin{itemize}[noitemsep]
  \item \emph{Почему природа выбрала именно $\mathrm{SiO_2}$ для хранения эталона времени?}
    Атом кремния несёт заряд $Z(\mathrm{Si}) = \mathbf{14 = N_\varphi + 1}$ ($13+1$), а пара атомов кислорода --- $Z(\mathrm{O}_2) = 2\times 8 = \mathbf{16 = 2|{\triangle}(3)|}$. Их сумма даёт молекулу кремнезёма:
    \[
      Z(\mathrm{SiO_2}) \;=\; 14 + 16 \;=\; \mathbf{30} \;=\; |\square(3)|\cdot d_{\mathrm{out}}(E) \quad (6\times 5 = 30),
    \]
    а двухканальная клетка $(\mu_0,\mu_1) = (14,16)$ лежит глубоко внутри устойчивого четверть-круга ($14^2 + 16^2 = 196 + 256 = \mathbf{452 \le 4\,096 = K_P^2}$) с собственной равновесной фазой $\varphi_{\mathrm{SiO_2}} = (30\times 41)\bmod 512 = \mathbf{206 = 316_8}$;
  \item \emph{Тригональная винтовая ячейка $(\mathrm{SiO_2})_3$ и точный порог Паули тетраэдра $\mathrm{SiO_4}$ ($Z = 46$):}
    в трёхмерном пространстве $d = 3$ кристалл $\alpha$-кварца скручивается в киральную (правую или левую) винтовую спираль из \textbf{ровно $d = 3$ формульных единиц $(\mathrm{SiO_2})_3$ на элементарную ячейку}, заряд которой равен:
    \[
      Z\bigl((\mathrm{SiO_2})_3\bigr) \;=\; 3\times 30 \;=\; \mathbf{90} \;=\; \hat{M}(3) - d_{\mathrm{out}}(V) \quad (97 - 7 = 90)!
    \]
    При этом каждый жёсткий кремнекислородный тетраэдр $\mathrm{SiO_4}$ (делимый по вершинам с соседями) имеет полный заряд вершин $Z(\mathrm{SiO_4}) = 14 + 4\times 8 = \mathbf{46}$ --- а число $\mathbf{46}$ является \textbf{точным минимальным целым порогом плотного спинорного перекрытия Паули} на решётке:
    \[
      46^2 \;=\; \mathbf{2\,116} \;>\; 2\,048 \;=\; \frac{K_P^2}{2} \;>\; \mathbf{2\,025} \;=\; 45^2!
    \]
    Именно поэтому при механическом сжатии тетраэдра $\mathrm{SiO_4}$ ($Z = 46$) мгновенно включается упругий спинорный отскок Паули (прямой пьезоэлектрический эффект: механическая деформация $\leftrightarrow$ фазовый импульс на кольце $\mathbb{Z}_{512}$!), а высокая добротность кристаллической решётки отсекает случайный тепловой шум ниже порога Гейзенберга $41$, превращая маленький кристалл $\mathrm{SiO_2}$ в идеальные механико-фазовые часы!
  \item \emph{Физика реального неидеального кристалла: почему главная примесь в кварце --- алюминий $\mathrm{Al}_{13}$ ($46\to 45$), как работает водородная очистка ($45+1=46$) и почему остаточные примеси не ломают прибор:}
    ни один природный или выращенный в автоклаве кристалл кварца не бывает абсолютно чистым. Главная примесь замещения в любом реальном кварце --- это соседний с кремнием ($Z=14$) атом \textbf{алюминия $Z(\mathrm{Al}) = \mathbf{13 = N_\varphi}$}. Посмотрим, что происходит в узле решётки, где $\mathrm{Si}_{14}$ заменён на $\mathrm{Al}_{13}$:
    \begin{enumerate}[label=(\alph*),noitemsep]
    \item В нескомпенсированном дефекте $[\mathrm{Al}_{13}\mathrm{O}_4]^-$ (дымчатый кварц / радиационный центр окраски) заряд тетраэдра падает ровно на единицу --- с $46$ до:
      \[
        Z\bigl([\mathrm{AlO_4}]^-\bigr) \;=\; 13 + 32 \;=\; \mathbf{45}, \qquad 45^2 \;=\; \mathbf{2\,025} \;<\; 2\,048 \;=\; \frac{K_P^2}{2}!
      \]
      Именно на примесном узле алюминия тетраэдр проваливается ровно на $1$ шаг ниже порога Паули ($46\to 45$), теряя упругий пьезоэлектрический отскок и рассеивая свет!
    \item Когда в промышленности проводят гидротермальную пассивацию и электроочистку кварца (\emph{sweeping}), выгоняя тяжёлые щелочные ионы натрия $\mathrm{Na}_{11}$ ($Z(\mathrm{AlO_4Na}) = 45 + 11 = \mathbf{56 = 7|{\triangle}|}$, дающие пик потерь) или лития $\mathrm{Li}_3$ ($Z(\mathrm{AlO_4Li}) = 45 + 3 = \mathbf{48 = 4|N|}$) и сажая рядом с алюминием протон водорода $\mathrm{H}^+$ ($Z = 1 = \mu_{\mathrm{vac}}$, дефект $\mathrm{Al{-}OH}$), полный заряд тетраэдра \textbf{в точности возвращается на порог Паули чистого кварца}:
      \[
        Z(\mathrm{AlO_4H}) \;=\; 13 + 32 + 1 \;=\; 45 + 1 \;=\; \mathbf{46} \;\equiv\; Z(\mathrm{SiO_4})!
      \]
    \item Наконец, почему даже оставшиеся после неидеальной очистки человека следы примесей ($\mathrm{Al}_{13}, \mathrm{Fe}_{26}, \mathrm{Na}_{11}$ на уровне $1\text{--}100\text{ ppm}$) не разрушают работу кварцевого резонатора? Потому что локальное рассеяние на редких точечных дефектах даёт вклад в 12-соседнюю голономию $\delta\varphi_{\mathrm{imp}} \le 40 < 41$, который \textbf{на $100\%$ обнуляется мёртвой зоной Гейзенберга} ($\hat{n}_E = 0$), а постоянный статический сдвиг собственной частоты реального кристалла ($\pm 20\text{ ppm}$ в даташите) целиком вычитается замкнутой петлей калибровки Типа~IV!
    \end{enumerate}
  \end{itemize}
\end{enumerate}

\section{Урок 27. Дискретная теория возмущений, конечная калибровка $(I + N_{\mathrm{pert}})^{-1} = I - N_{\mathrm{pert}} + N_{\mathrm{pert}}^2$ и развенчание мифов о «10/11/26 измерениях струн» и «Мультивселенных»}

В классической университетской физике любая реальная задача с примесями, трением и неидеальными деталями решается через \textbf{теорию возмущений}: к идеальному оператору $H_0$ добавляют малое возмущение $\lambda V$ и раскладывают ответ в бесконечный ряд $\Psi^{(0)} + \Psi^{(1)} + \Psi^{(2)} + \dots + \Psi^{(\infty)}$. Но в непрерывном континууме $\mathbb{R}$ этот бесконечный ряд почти всегда \emph{расходится}, порождая искусственные бесконечности и заставляя теоретиков выдумывать «перенормировки», «10 или 11 измерений теории струн» и «ветвящиеся Мультивселенные».

Посмотрим, как устроена теория возмущений и калибровки на нашем дискретном 9-битном кольце $\mathbb{Z}_{512}$ и кубооктаэдре $O_2(3)$:

\begin{enumerate}[label=\textbf{\arabic*.},leftmargin=*]
\item \textbf{Почему на $\mathbb{Z}_{512}$ ряд теории возмущений конечен и никогда не расходится:}
  \begin{itemize}[noitemsep]
  \item \emph{Нулевой порядок (Мёртвая зона Гейзенберга $|\delta\varphi|\le 40 < 41$ отрезает бесконечный хвост в точный ноль):}
    в непрерывном анализе даже микроскопическая помеха $10^{-50}$ тянет за собой бесконечный хвост производных Тейлора. На кольце $\mathbb{Z}_{512}$ любое возмущение внутри коридора $[-40,+40]$ ($81 = 3^4$ состояний) даёт $\hat{n}_E(\delta\varphi) = \mathbf{0}$ --- малый шум обрезается природой в \textbf{точный ноль};
  \item \emph{Максимум $6 = |\square(3)|$ ненулевых квантовых порядков возмущения на ячейке:}
    поскольку потолок Паули равен $\Phi_{\mathrm{Pauli}} = 256$, а квант порога Гейзенберга равен $\Delta\varphi_{\mathrm{disc}} = 41$, максимальный физически возможный порядок энергетического возмущения на одной ячейке равен:
    \[
      n_{\max} \;=\; \left\lfloor \frac{\Phi_{\mathrm{Pauli}}}{\Delta\varphi_{\mathrm{disc}}} \right\rfloor \;=\; \left\lfloor \frac{256}{41} \right\rfloor \;=\; \mathbf{6 \;=\; |\square(3)|},
      \qquad
      256 - 6\times 41 \;=\; \mathbf{10 \;=\; Z(\mathrm{Ne})}!
    \]
    Ровно по одному квантовому порядку на каждую из $6$~квадратных граней кубооктаэдра --- никакого 7-го, 100-го или бесконечного порядка на ячейке физически не существует!
  \item \emph{Точный обрыв ряда Неймана на 2-м порядке для калибровки прибора $(I + N_{\mathrm{pert}})^{-1} = I - N_{\mathrm{pert}} + N_{\mathrm{pert}}^2$:}
    пусть межканальная наводка и примеси в приборе задаются строго треугольным иерархическим оператором возмущения $N_{\mathrm{pert}}$ по $d = 3$ пространственным слоям ($N_{\mathrm{pert}}^3 \equiv 0$). Тогда бесконечный ряд теории возмущений Неймана $(I + N_{\mathrm{pert}})^{-1} = \sum_{k=0}^\infty (-1)^k N_{\mathrm{pert}}^k$ \textbf{обрывается навсегда на втором порядке}:
    \[
      \mathcal{K}_{\mathrm{cal}} \;=\; (I + N_{\mathrm{pert}})^{-1} \;\equiv\; \mathbf{I \;-\; N_{\mathrm{pert}} \;+\; N_{\mathrm{pert}}^2},
    \]
    восстанавливая невозмущённый сигнал прибора $(s_0,s_1,s_2)$ \textbf{без единого деления и с нулевой ошибкой усечения}!
  \end{itemize}
\item \textbf{Откуда в теории струн и М-теории взялись числа $10$, $11$ и $26$ (и почему в пространстве всегда ровно $d = 3$ измерения!):}
  поскольку диофантово уравнение Каталана $1 + 2^d = d^2$ имеет единственное решение $d = 3$, внешнее физическое пространство строго \textbf{трёхмерно}. Почему же в формулах теории струн при сокращении аномалий всплывали числа $10$, $11$ и $26$? Потому что теоретики приняли \emph{внутренние 9-битные дорожки памяти одной 3D-клетки} и \emph{26 векторов кубической оболочки} за «дополнительные оси пространства»:
  \begin{itemize}[noitemsep]
  \item \textbf{Число $10$ (суперструны):}
    \[
      d + d_{\mathrm{out}}(V) \;=\; 3 + (2^3 - 1) \;=\; 3 + 7 \;=\; \mathbf{10},
      \qquad\text{или}\qquad
      d_{\mathrm{in}}(\square) + |\square(3)| \;=\; 4 + 6 \;=\; \mathbf{10}
    \]
    ($3$ внешние оси пространства $+$ $7$ независимых относительных дорожек двухтактового регистра клетки, либо $4$ координаты пространства-времени $+$ $6$ квадратных граней кубооктаэдра!);
  \item \textbf{Число $11$ (М-теория):}
    \[
      d + |{\triangle}(3)| \;=\; 3 + 2^3 \;=\; 3 + 8 \;=\; \mathbf{11}
    \]
    ($3$ внешние пространственные координаты $(x,y,z)$ $+$ все $8$~девятибитных кольцевых дорожек \texttt{Ring9} в 72-битном двухтактовом регистре клетки \texttt{VerletCell72}, $B_2(3) = 8\times 9 = 72\text{ бита}$!);
  \item \textbf{Число $26$ (бозонная струна):}
    \[
      |\square(3)| + |{\triangle}(3)| + |N_{12}(3)| \;=\; 6 + 8 + 12 \;=\; \mathbf{26 \;=\; 3^3 - 1}
    \]
    (все $26$ ненулевых векторов элементарного куба Мура $3\times 3\times 3$ вокруг одной 3D-клетки: $6$~граней $O_1$, $12$~рёбер кубооктаэдра $O_2$ и $8$~углов $O_3$!).
  \end{itemize}
\item \textbf{Почему не существует никаких «параллельных Мультивселенных» ($64\times 81 = 5\,184 = 72^2$ скрытых микросостояний ОДНОЙ Вселенной):}
  грубый макроскопический Т-прибор не видит $K_P(3) = 64$ градации мантиссы Вейля--Швингера и $|\mathcal{E}_{\mathrm{Heis}}(3)| = 81 = 3^4$ подпороговое состояние мёртвой зоны Гейзенберга $[-40,+40]$, склеивая на каждой ячейке ровно:
  \[
    K_P(3)\cdot |\mathcal{E}_{\mathrm{Heis}}(3)| \;=\; 64 \times 81 \;=\; \mathbf{5\,184 \;=\; 72^2 \;=\; \bigl(B_2(3)\bigr)^2}
  \]
  различных детерминированных микросостояний М-слоя в одно макро-наблюдение. То, что в XX~веке приняли за «расщепление мира на параллельные Мультивселенные при каждом квантовом измерении», есть просто движение этих $72^2 = 5\,184$ подпороговых микросостояний внутри \textbf{одной-единственной детерминированной 3D-Вселенной}, растущей из зерна $\mathrm{IV}(3) = 263 = 407_8$!
\item \textbf{Анатомия второго узла стенда: силовой затвор MOSFET ($\mathrm{Si + SiO_2}$), медная катушка ($Z(\mathrm{Cu}) = 29$) и прозрачный мост «Кварцевая стенка чашки Петри $\leftrightarrow$ Воздух» ($30 \equiv 30$):}
  проследим теперь, что происходит с тактовым импульсом после того, как он вышел из кварцевого резонатора $\alpha$-$\mathrm{SiO_2}$ ($27\text{ МГц}$, период кольца $512$ тактов $f_{512} = 52\,734{,}375\text{ Гц}$, $1\text{ тик} = 37{,}037\text{ нс}$) и пошёл через силовой ключ MOSFET и медную катушку в чашку Петри:
  \begin{itemize}[noitemsep]
  \item \emph{Структура затвора MOSFET ($\mathrm{Si} + \mathrm{SiO_2}$, $Z = 14 + 30 = 44$):}
    сам полевой транзистор MOSFET --- это слоистый конденсатор «кремниевый канал $\mathrm{Si}_{14}$ $+$ подзатворный диэлектрик $\mathrm{SiO_2}_{30}$» с суммарной оболочкой $Z(\mathrm{Si}) + Z(\mathrm{SiO_2}) = 14 + 30 = \mathbf{44 = Z(\mathrm{Ru})}$. При перезаряде входной ёмкости затвора $C_{\mathrm{iss}}$ и прохождении плато Миллера ключ вносит задержку фронта $\delta\varphi_{\mathrm{gate}}$ (например, $\sim 925\text{ нс} \approx \mathbf{25 = d_{\mathrm{out}}(E)^2}$ тиков кольца);
  \item \emph{Медная катушка индуктивности ($Z(\mathrm{Cu}) = 29 = \mathrm{magic}(28) + 1$) и её реактивный фазовый сдвиг:}
    медный провод не случайно обладает рекордной проводимостью: атом меди имеет заряд $Z(\mathrm{Cu}) = 28 + 1 = \mathbf{29}$ --- ровно один свободный электрон $\mu_{\mathrm{vac}} = 1$ поверх полностью замкнутой первой магической оболочки никеля $Z = 28$! На частоте кольца $f_{512} = 52{,}73\text{ кГц}$ индуктивность катушки $L$ отстаёт по току от напряжения на угол $\mathrm{tan}_3(\delta\varphi_L) = (\omega L,\,R)$. В частности, при $\omega L = R$ катушка задерживает фазу тока ровно на четверть квадранта ($45^\circ$), то есть на $\delta\varphi_L = \mathbf{64 = K_P(3)}$ тика!
  \item \emph{Акустическое и диэлектрическое равенство оболочек на границе «Стекло чашки Петри $\leftrightarrow$ Комнатный воздух» ($30 \equiv 30$):}
    почему стеклянная стенка чашки Петри ($\mathrm{SiO_2}$) не блокирует передачу фазового поля в воздушную полость камеры? Потому что заряд молекулы кварцевого стекла стенки:
    \[
      Z(\mathrm{SiO_2}) \;=\; Z(\mathrm{Si}) + Z(\mathrm{O}_2) \;=\; 14 + 16 \;=\; \mathbf{30}
    \]
    \textbf{тождественно совпадает} с суммой зарядов двухатомных молекул азота и кислорода комнатного воздуха:
    \[
      Z(\mathrm{N}_2) + Z(\mathrm{O}_2) \;=\; 2\times 7 + 2\times 8 \;=\; 14 + 16 \;=\; \mathbf{30}!
    \]
    И атом кремния $\mathrm{Si}$ ($Z=14$), и молекула азота $\mathrm{N}_2$ ($Z=14$) несут один и тот же заряд $14 = N_\varphi + 1$, спаренный с тем же самым кислородом $\mathrm{O}_2$ ($Z=16$)!
  \item \emph{Точный расчёт предкомпенсации прошивки FPGA ещё до включения питания:}
    зная по даташиту задержку затвора $\delta\varphi_{\mathrm{gate}}$ и индуктивный сдвиг катушки $\delta\varphi_L$, мы задаём в регистре \texttt{Tang Nano 9K} упреждающую фазу:
    \[
      \varphi_{\mathrm{FPGA}} \;=\; \bigl(A_{\mathrm{coh}} + \delta\varphi_L + \delta\varphi_{\mathrm{gate}}\bigr)\bmod 512
      \;=\; (139 + 64 + 25)\bmod 512 \;=\; \mathbf{228 \;=\; 344_8}.
    \]
    Тогда, пройдя через задержку затвора ($-25$ тиков) и индуктивность медной катушки ($-64$ тика), магнитный и акустический фронт входит сквозь согласованную стеклянную стенку ($30 \equiv 30$) внутрь чашки Петри \textbf{ровно на фазе каталитического резонанса}:
    \[
      \varphi_{\mathrm{dish}} \;=\; (228 - 64 - 25)\bmod 512 \;=\; \mathbf{139 \;=\; A_{\mathrm{coh}}(3)} \;\implies\; D\bigl(K_P, 139\bigr) \;=\; \mathbf{\varepsilon \;=\; 8\,192}!
    \]
  \end{itemize}
\end{enumerate}

\section{Урок 28. Единый ньютоновский гомеоморфизм всех шести отраслей инженерии, четырёхслойный кремниевый гетеропереход и сквозная автокалибровка 12-лучевого резонатора кубооктаэдра}

В завершающем уроке курса сведём воедино все шесть отраслей прикладной инженерии и покажем, как на основе трёх законов Ньютона в целых числах и пяти порогов кольца $\mathbb{Z}_{512}$ ($|\Delta\varphi|\le 40 < 41$, $K_P = 64$, $Q_\square = 128$, $\Phi_{\mathrm{Pauli}} = 256$) проектируются и калибруются реальные установки --- от четырёхслойного кремниевого гетероперехода до 12-катушечного резонатора кубооктаэдра на столе.

\begin{table}[htbp]
\centering
\caption{Шесть отраслей инженерии: континуальные ограничения $\mathcal{T}$-слоя против точной ньютоновской фазовой механики на $O_2(3)\times\mathbb{Z}_{512}$.}
\label{tab:course-six-engineering-branches}
\footnotesize
\begin{tabular}{@{}p{0.21\textwidth}p{0.35\textwidth}p{0.39\textwidth}@{}}
\toprule
\textbf{Отрасль инженерии} & \textbf{Проблема в континуальной $\mathcal{T}$-модели ($\mathbb{R}^3$)} & \textbf{Решение в ньютоновской $\mathcal{M}$-механике ($O_2(3)\times\mathbb{Z}_{512}$)} \\
\midrule
\textbf{1. Электротехника и силовая электроника} &
Гармоники, реактивная мощность, пульсации конденсаторов и джоулевы потери $I^2 R$. &
$V=\mathcal{C}_3,\ I=\mathcal{S}_3$ на $L^1$-ромбе: $|V|+|I|\equiv 128$ ($0\%$ пульсаций без конденсатора), MPPT при $k^*=64$ ($P_{\max}=4\,096$), гашение взаимной индукции 12 катушек $4(+1)+2(0)+4(-1)=0$ и $\hat{n}_E=0$ при $|\Delta\varphi|\le 40$. \\
\addlinespace[3pt]
\textbf{2. Полупроводники и фотовольтаика} &
Термализация высокоэнергетических экситонов в тепло и рекомбинация на интерфейсе. &
Бесфононный шаг шва $\lfloor 256/12\rfloor = 21 < 41$ ($\hat{n}_E=0$); кислородный замок седловины $Z(\mathrm{C_{18}H_{12}})+Z(\mathrm{O})=120+8=128=Q_\square$, барьер пассивации Паули $\mathrm{SiO}_4$ $45^2 < 2\,048 < 46^2 = 2\,116$ и окно Френеля $288/289$. \\
\addlinespace[3pt]
\textbf{3. Вычислительная техника и защита} &
Тепло необратимых булевых вентилей $\ge k_{\mathrm{B}}T\ln 2$ и уязвимость гипотез сложности. &
Обратимое АЛУ Верле Канала~1 (\texttt{Ring9/Bus36/VerletCell72}) с четверть-квадратным умножением $(x+y)^2-(x-y)^2=4xy$ при $\hat{n}_E=0$; барьеры решателя внутри носителя ($72^2=5\,184$, $7^T$) и встроенный шифр Вернама $H(M\mid T)=H(M)$. \\
\addlinespace[3pt]
\textbf{4. Теплотехника и термодинамика} &
Трактует тепло как необратимый хаос, требующий внешнего холодильника или стирания памяти. &
$12$-лучевой расщепитель кубооктаэдра без памяти делит даже $6$-квантовый удар Паули $\Phi_{\mathrm{Pauli}}=256$ на $\lfloor 256/12\rfloor = 21 < 41$, переводя тепловой удар в баллистический ток ($12\times 21=252>0$) при $\hat{n}_E(21)=0$. \\
\addlinespace[3pt]
\textbf{5. Гидрогазодинамика и акустика} &
Уравнения Навье--Стокса дают континуальные сингулярности и неквантованную турбулентность. &
$4$-канальный нильпотентный каскад $\mathcal{M}_{\mathrm{chan}} = I_4 + N_{\mathrm{nil}}$ ($N_{\mathrm{nil}}^3=0$, $\det(2\mathcal{M})=24$) точно обращается полиномом $I_4 - N_{\mathrm{nil}} + N_{\mathrm{nil}}^2$; турбулентность квантуется вихрями $512 n_f$ выше порога $41$, критическое число Маха $c_0/c=\sqrt{2}$. \\
\addlinespace[3pt]
\textbf{6. Сопромат, катализ и ТММ приводов} &
Эмпирические коэффициенты усталости, накопление ошибки Рунге--Кутты и фазовый джиттер каскадов. &
Упругий режим без усталости при $|\Delta\varphi|\le 40 < 41$; седловина катализа в верхней мёртвой точке $Q_\square=128$ (коллапс барьера в $587$ раз); потолок Паули $D_{\max}=833\varepsilon$ при $256$; точный привод Верле $\Delta^2\varphi\equiv\Phi_{\mathrm{imp}}$ и гашение джиттера $(\varphi_{\mathrm{target}}-\Sigma j)+\Sigma j = \varphi_{\mathrm{target}}$. \\
\bottomrule
\end{tabular}
\end{table}

\begin{enumerate}[label=\textbf{\arabic*.},leftmargin=*]
\item \textbf{Четырёхслойный бесфононный кремниевый гетеропереход ($120 + 8 = 128 = Q_\square(3)$ и $45^2 < 2\,048 < 46^2$):}
  почему обычный кремниевый переход ($\mathrm{Si}$, $Z=14$) теряет большую часть избыточной энергии коротковолнового фотона в паразитное тепло? Потому что без согласования оболочек скачок фазы на интерфейсе пробивает порог Гейзенберга $41$! В четырёхслойной гетероструктуре:
  \begin{itemize}[noitemsep]
  \item \emph{Слой~1 (Кварцевый микрорельеф $O_2(3)$):} отношение показателей преломления $n_1/n_2 = \mathrm{mod\_bits}/|{\triangle}| = 9/8$ даёт просветление Френеля $T_{\mathrm{Fresnel}} = 1 - ((9-8)/(9+8))^2 = \mathbf{288/289\approx 99{,}654\%}$;
  \item \emph{Слой~2 (Кристаллический тетрацен $\mathrm{C_{18}H_{12}}$, $Z = 18\times 6 + 12 = \mathbf{120 = 128 - 8}$):} делит один синглетный экситон ($S_1 \to T_1 + T_1$) на пару триплетных экситонов при распределении полупериодного импульса Паули по $12$ лучам кубооктаэдра с шагом $\Delta\varphi_{\mathrm{seam}} = \lfloor 256/12\rfloor = 21 < 41$;
  \item \emph{Слой~3 (Кислородно-силанольный седловой мост $\mathrm{SiO_4}/[\mathrm{AlO_4H}]$, $Z = 46$):} мостиковый атом кислорода $Z(\mathrm{O}) = 8$ доводит оболочку тетрацена точно до верхней мёртвой точки ротатора $120 + 8 = \mathbf{128 = Q_\square(3)}$ ($\mathcal{C}_3(128)=0,\ \mathcal{S}_3(128)=+128$), снимая обменный барьер $12^2 = 144\text{ мэВ}$ на голом $\mathrm{Si}(100)$, а тетраэдр $Z(\mathrm{SiO_4}) = Z(\mathrm{AlO_4H}) = 46$ ($46^2 = 2\,116 > 2\,048 = K_P^2/2 > 2\,025 = 45^2 = Z(\mathrm{AlO_4})^2$) запирает утечку неразделённого синглета;
  \item \emph{Слой~4 (Базовый переход IBC-$\mathrm{Si}$ с тыльными контактами):} принимает триплетную пару при шаге переноса $\Delta\varphi_{\mathrm{seam}} = 21 < 41$ по первому закону Ньютона в зазоре зацепления с \textbf{нулевым рождением тепловых фононов ($\hat{n}_E(21) = 0$)}!
  \end{itemize}
\item \textbf{Диагонализация взаимной индукции 12 катушек кубооктаэдра $O_2(3)$ и сквозная автокалибровка стенда:}
  при сборке лабораторного 12-лучевого резонатора вокруг кварцевой кюветы ($\mathrm{SiO_2}$, $Z = 14 + 16 = \mathbf{30} \equiv Z(\mathrm{N_2}) + Z(\mathrm{O_2})$):
  \begin{itemize}[noitemsep]
  \item для любой выбранной катушки $e_0\in O_2(3)$ ($|e_0|^2 = 2$) остальные $11$ катушек распадаются на $4$ соседа под $60^\circ$ ($e_0\cdot e_j = +1$), $2$ ортогональных соседа под $90^\circ$ ($e_0\cdot e_j = 0$), $4$ тупоугольных соседа под $120^\circ$ ($e_0\cdot e_j = -1$) и $1$ антиподального партнёра под $180^\circ$ ($-e_0$), поэтому сумма внеосевых коэффициентов взаимной индукции тождественно обращается в ноль:
    \[
      4\times(+1) \;+\; 2\times 0 \;+\; 4\times(-1) \;=\; \mathbf{0};
    \]
  \item любые аппаратные задержки реального тракта произвольной длины $\texttt{jitters} = [\delta\varphi_1,\dots,\delta\varphi_m]$ (задержка затвора MOSFET $\delta\varphi_{\mathrm{gate}} = 25$ тиков, индуктивный сдвиг медной катушки $Z(\mathrm{Cu})=29$ $\delta\varphi_L = 64$ тика, стенка кюветы $\delta\varphi_{\mathrm{wall}}$) и межканальная наводка 4 датчиков $(2I_4 + N_{\mathrm{nil}})$ точно обращаются за один такт замкнутой петли ФАПЧ:
    \[
      \bigl(\varphi_{\mathrm{target}} - \Sigma\texttt{jitters}\bigr) + \Sigma\texttt{jitters} \;\equiv\; \mathbf{\varphi_{\mathrm{target}}},
      \qquad
      (12 I_4 - 6 N_{\mathrm{nil}} + 3 N_{\mathrm{nil}}^2)(2 I_4 + N_{\mathrm{nil}})\mathbf{s} \;\equiv\; \mathbf{24\,\mathbf{s}}\quad\forall\,\mathbf{s}\in\mathbb{Z}^4!
    \]
  \end{itemize}
\end{enumerate}

